\section{Online algorithm and state accounting}
\label{app:algorithm}
\subsection{Initialization, routing, and state lifetime}
The two residual matrices and all candidate scores start at zero. The stored
counter starts at zero. At the start of each episode, the current route is
source, the route is unlocked, and the temporary signature and pending-loss
lists are cleared. The matrices and cumulative scores persist across episodes
within a run. A new run creates a new state.

The four routing classes are source, shear, attenuation, and noise. The router
is a frozen two-layer network with a 64-unit ReLU hidden layer, fit offline
using calibration signatures. Online, it receives the mean signature of the
first four transitions. A signature contains candidate losses relative to identity, signed and
absolute identity prediction errors, signed and absolute commands, and current
physical observations. It also contains signed and absolute observation changes
and cross-products between the first two command coordinates and the first
six observation-change coordinates. These concatenate to 247
values per transition. The true mechanism label is used to fit the router
offline but is absent from its online input.

Once the fourth signature arrives, the router locks for that episode. For
shear or attenuation, the four pending loss vectors are added to the persistent
scores, and their count is committed. Subsequent outcomes add one new loss
vector each. Source and noise routes keep the identity command transformation
and do not update residual matrices. Both temporary lists are cleared when
the route locks. The scores are shared across tasks and indexed by candidate;
the two matrices are indexed by mechanism family.

The selected candidate is the minimum cumulative score within the locked
family. Dividing scores by the shared count would leave this ordering
unchanged. Ties use the first minimum in the fixed bank order, matching the
implementation's \texttt{argmin}. Scores are computed from the frozen
predictor, before the residual-matrix correction.

\subsection{Candidate transformations and command correction}
The identity candidate is followed by 15 shear values
$k\in\{0.05,0.10,\ldots,0.75\}$ and 16 attenuation values
$g\in\{0.20,0.25,\ldots,0.95\}$. For command coordinates $(a_x,a_y)$,
\begin{align}
 T_{\rm shear}(a;k)_x&=a_x,&
 T_{\rm shear}(a;k)_y&=a_y+ka_x,\\
 T_{\rm attenuation}(a;g)_{x,y}&=g\,a_{x,y}.&&
\end{align}
Other coordinates are unchanged. Transformed candidates are clipped to
$[-1,1]$ before querying the predictor. The inferred inverse uses
$\tilde a_y=a_y^*-\hat k a_x^*$ for shear, or
$\tilde a_{x,y}=a_{x,y}^*/\max(\hat g,0.20)$ for attenuation.
The hat denotes the selected bank value. The floor $0.20$ prevents excessive
division and coincides with the lowest attenuation candidate. The change
from the desired command is then capped by Equation~\eqref{eq:command}.
If the resulting command exceeds the action bounds, the entire proposal is
replaced by the desired command for that example.

\subsection{Exact matrix write}
\label{app:matrix-update}
The fixed basis $B$ contains the first 30 right singular vectors of centered
offline action hidden features. After observing $d_n$ and committing the
candidate-score update, the prediction used for the matrix write is recomputed
with the selected candidate and the current matrix. Let $C$ denote this
active matrix, and let
\begin{equation}
 L_n(C)=\frac1{D_q}\|W_{\rm out}(h_n+Cz_n)+b_{\rm out}-d_n\|_2^2,
 \qquad A_n=\nabla_C L_n(C).
\end{equation}
Here $L_n$ is the current matrix-write loss, $A_n$ is its matrix gradient, and
all other quantities are held fixed
for this write. The normalized step and projection are
\begin{equation}
 C'=C-\eta\frac{A_n}{\max(\|A_n\|_F,\epsilon)},\qquad
 C\leftarrow C'\min\!\left(1,\frac{R_C}{\max(\|C'\|_F,\epsilon)}\right),
 \quad \epsilon=10^{-12}.
 \label{eq:update}
\end{equation}
$C'$ is the tentative matrix before projection, and $\epsilon$ safeguards
division by a zero norm. The Frobenius norm $\|\cdot\|_F$ is the square root of the sum of squared
matrix entries. Normalization controls step length; projection bounds the
matrix norm. A zero gradient yields a zero step. No optimizer moments are
maintained. This recomputed training loss is distinct from the pre-update
prediction error already recorded for evaluation.

\paragraph{A useful interpretation.}
Let $e_n=d_n-g_{\theta_0}(q_n,T_j(a_n),c)$ be the uncorrected error after
candidate selection, and let $v_n=W_{\rm out}C_\rho z_n$ be the correction in
output space. Expanding the square gives
\begin{equation}
 \frac{\|e_n\|_2^2-\|e_n-v_n\|_2^2}{D_q}
 =\frac{2e_n^\top v_n-\|v_n\|_2^2}{D_q}.
\end{equation}
The first term rewards alignment with the residual; the second penalizes an
oversized correction. Thus prediction improvement depends on its direction
as well as its magnitude. This identity concerns prediction error, while
the separate inverse determines command changes.

\subsection{One decision cycle}
\begin{enumerate}[leftmargin=*,itemsep=3pt]
\item Read the current physical observation and desired policy command.
      Use the previously available route and scores to form a capped command.
\item Predict its outcome with the frozen core and the currently available
      residual matrix. Save this prediction for evaluation.
\item Execute the command and wait for the outcome. In live attenuation
      control, the actuator-transformed command is held over the macro horizon.
\item Score the saved prediction against the observed outcome before writing
      either external component.
\item Compute all candidate losses with the uncorrected frozen predictor.
      Before route lock, append the bounded signature/loss prefix. At the
      fourth outcome, infer the family and, for shear or attenuation, commit
      the four pending loss vectors exactly once. Clear the temporary prefix.
      On subsequent outcomes in these routes, add only the new loss vector.
\item If the route is now locked to shear or attenuation, perform one matrix
      write with Equation~\eqref{eq:update}. The fourth outcome is the first
      eligible write; its candidate losses are not added a second time.
\item Continue with the next decision. The first opportunity for a changed
      command is decision five.
\end{enumerate}
The code provides separate \texttt{predict}, \texttt{inverse}, and
\texttt{observe} operations to keep these read/write times explicit.

\subsection{Complete mutable-state ledger}
\begin{table}[tbp]\centering\small
\caption{Peak mutable numeric allocation in FP32. Fixed model/router/basis
parameters and transient activations are outside this adaptation-state metric.}
\label{tab:ledger}
\begin{tabular}{@{}llr@{}}\toprule
Object & Arithmetic / lifetime & Bytes\\\midrule
Two residual matrices & $2\times96\times30\times4$; stream & 23,040\\
Candidate scores and count & $(32+1)\times4$; stream & 132\\
Compatibility padding & $303\times4$; fixed allocation & 1,212\\
Routing signatures and losses & $4\times(247+32)\times4$; prefix & 4,464\\\midrule
Peak total & Tensor block plus prefix & 28,848\\\bottomrule
\end{tabular}\end{table}
The allocated tensor block is 24,384 bytes. The active entries use 23,172
bytes; the remaining entries preserve the historical allocation layout.
The prefix is bounded at four observations and released on route lock.
Python-object overhead and transient inference/gradient buffers are not
included in this numeric-state ledger. The same persistent-state accounting
includes trainable prediction parameters, optimizer tensors, and replay for
the comparison arms. Bytes are not a hardware energy measurement.
